\chapter{Implementation}

This chapter provides additional implementation details, development practices, and technical decisions beyond the architecture described in Chapter 4.

\section{Technology Stack}

\begin{description}
    \item[Language:] Python 3.12 for all application code, analysis scripts, and tests. Chosen for AWS Lambda runtime support, rich scientific computing ecosystem (NumPy, Pandas, SciPy), and rapid prototyping.
    
    \item[Infrastructure as Code:] AWS SAM (Serverless Application Model) for CloudFormation template generation. Preferred over raw CloudFormation for concise Lambda + event source mapping syntax.
    
    \item[Testing:] pytest with moto (mocked AWS services), cfn-lint for template validation, GitHub Actions for CI/CD.
    
    \item[Analysis:] NumPy, Pandas, SciPy (Mann-Whitney U, Kruskal-Wallis H, bootstrap), Matplotlib, Seaborn for visualization.
    
    \item[Version Control:] Git with semantic commit messages, branch protection, automated testing before merge.
\end{description}

\section{Development Workflow}

\subsection{Incremental Build}

1. \textbf{SAM Template + Handlers:} Implemented Lambda handlers first with local moto tests.\\
2. \textbf{Fault Injection:} Added fault switch as pluggable module, validated with unit tests.\\
3. \textbf{Simulator:} Built discrete-event engine, validated SQS semantics against AWS documentation.\\
4. \textbf{Experiment Runner:} Config matrix generator, manifest persistence, seeded randomization.\\
5. \textbf{Analysis:} Hypothesis tests, effect sizes, corrections, figures.\\
6. \textbf{Documentation:} Config manual, architecture diagram, demo walkthrough, viva Q\&A.

\subsection{Test-Driven Approach}

Each component was developed with accompanying unit tests before integration:
\begin{itemize}
    \item Fault injection: Test each fault type in isolation.
    \item SQS simulator: Test visibility timeout expiry, receive count tracking, DLQ redrive.
    \item Metrics: Test loss, duplicate, recovery time calculations with known inputs.
    \item Idempotency: Test deduplication logic with repeated \texttt{order\_id}.
\end{itemize}

Integration tests validated end-to-end flow:
\begin{itemize}
    \item Deploy SAM stack with moto, send messages, verify Lambda invocation.
    \item Run pilot campaign in simulator, verify manifests contain expected metrics.
\end{itemize}

\section{Simulator Validation}

The simulator's fidelity was validated through:

\subsection{Unit Test Coverage}

\begin{itemize}
    \item \texttt{test\_sqs\_sim.py}: 15 test cases covering send, receive, delete, visibility timeout, DLQ redrive, receive count.
    \item \texttt{test\_engine.py}: 8 test cases for event ordering, time advancement, concurrent events.
    \item \texttt{test\_handler.py}: 10 test cases for batch processing, partial failures, idempotency.
\end{itemize}

\subsection{Manual Verification}

\begin{enumerate}
    \item \textbf{Visibility Timeout:} Sent 1 message, received it, waited VT seconds without delete → verified message reappeared with incremented receive count.
    \item \textbf{DLQ Redrive:} Set \texttt{maxReceiveCount=3}, injected unhandled error → verified message moved to DLQ after 3 receives.
    \item \textbf{Partial Failure:} Batch of 10 messages, 5 succeed, 5 fail → verified \texttt{batchItemFailures} causes only 5 to be redelivered.
\end{enumerate}

\subsection{Comparison with AWS Documentation}

Simulator behavior matches AWS SQS Developer Guide~\citep{AWSSQSGuide}:
\begin{itemize}
    \item At-least-once delivery: Messages may be delivered multiple times (even without faults, due to eventual consistency simulation).
    \item Visibility timeout: Message invisible for configured duration; reappears if not deleted.
    \item \texttt{maxReceiveCount}: After N receives without delete, message moves to DLQ.
    \item Batch processing: Lambda receives up to \texttt{BatchSize} messages; \texttt{ReportBatchItemFailures} allows partial success.
\end{itemize}

\section{Code Quality and Maintainability}

\subsection{Linting and Formatting}

\texttt{ruff} (fast Python linter) enforces PEP 8, detects unused imports, checks type hints. \texttt{cfn-lint} validates SAM template against CloudFormation spec. CI/CD fails on lint errors.

\subsection{Documentation}

\begin{itemize}
    \item Docstrings for all public functions (Google style).
    \item \texttt{README.md} with quick start, layout, commands.
    \item \texttt{docs/CONFIGURATION\_MANUAL.md}: Step-by-step guide for live AWS deployment.
    \item \texttt{docs/ARCHITECTURE.md}: System diagram, data flow, fault injection design.
    \item \texttt{docs/DEMO\_WALKTHROUGH.md}: Script for demo video.
    \item \texttt{docs/VIVA\_QA.md}: Anticipated viva questions and answers.
\end{itemize}

\subsection{Dependency Management}

\texttt{requirements.txt} pins versions (e.g., \texttt{boto3>=1.34}) for reproducibility. \texttt{requirements-dev.txt} adds test and lint dependencies. \texttt{scripts/scaffold.sh} creates virtual environment and installs dependencies.

\section{Cost Management}

\subsection{Cost Estimation}

\texttt{scripts/estimate\_cost.py} computes expected AWS charges for a config:
\begin{itemize}
    \item SQS requests: \$0.40 per million (beyond free tier).
    \item Lambda invocations: \$0.20 per million + \$0.0000166667 per GB-second.
    \item DynamoDB writes: \$1.25 per million.
\end{itemize}

Pilot (18 runs, 1000 messages each) estimated at \$0.10. Full experiments (350+ runs, 10,000 messages each) estimated at \$50-100. Simulator avoids all costs.

\subsection{Free-Tier Guard}

\texttt{scripts/assert\_free\_tier\_guard.py} checks if a config exceeds AWS Free Tier limits (1M SQS requests, 1M Lambda invocations per month). Exits with error if exceeded, preventing accidental charges during development.

\section{Makefile Targets}

\texttt{Makefile} provides shortcuts for common tasks:
\begin{itemize}
    \item \texttt{make setup}: Create venv, install dependencies.
    \item \texttt{make test}: Run pytest (unit + integration).
    \item \texttt{make lint}: Run ruff on src/, tests/, analysis/.
    \item \texttt{make validate}: Run cfn-lint on template.yaml.
    \item \texttt{make pilot}: Execute 18 pilot runs (simulator).
    \item \texttt{make baseline}: Steady-state replication (no faults).
    \item \texttt{make arms}: Sync vs. queue comparison.
    \item \texttt{make campaigns}: Fault campaigns A-E.
    \item \texttt{make burst}: Burst load replication.
    \item \texttt{make experiments}: All of the above (pilot + baseline + arms + campaigns + burst).
    \item \texttt{make stats}: Run hypothesis tests, generate summary tables.
    \item \texttt{make figures}: Generate 11 plots from results.
    \item \texttt{make gates}: Quality gates (tests + validate + self-check + figures exist).
    \item \texttt{make deploy}: Deploy SAM stack to AWS (DRY\_RUN=0 required).
    \item \texttt{make destroy}: Delete SAM stack from AWS.
\end{itemize}

All targets default to \texttt{DRY\_RUN=1} (simulator). Override with \texttt{DRY\_RUN=0 make <target>} for live AWS.

\section{Challenges and Solutions}

\subsection{Challenge: Simulating At-Least-Once Delivery}

SQS Standard queues may deliver messages more than once even without failures, due to distributed architecture. Modeling this required adding a small probability of duplicate delivery (~0.2\%) in the simulator, validated against AWS-observed baseline from \citet{Kyrychenko2025}.

\textbf{Solution:} Added \texttt{duplicate\_send\_prob} parameter to simulator config; randomly reinserts \texttt{MessageVisible} event with low probability.

\subsection{Challenge: Partial Failure Reporting Semantics}

Lambda's \texttt{ReportBatchItemFailures} requires returning specific \texttt{itemIdentifier} (SQS message ID). Simulator needed to track which messages failed and only reinsert those.

\textbf{Solution:} Handler returns list of failed \texttt{order\_id}, simulator cross-references against message IDs, deletes successes, leaves failures in-flight.

\subsection{Challenge: Recovery Time Definition}

Multiple definitions possible: (1) last message processed, (2) backlog returns to baseline, (3) visible-only depth returns. Choice affects interpretation.

\textbf{Solution:} Report both (2) and (3); \texttt{results/README.md} documents definitions; analysis interprets both.

\section{Lessons Learned}

\begin{itemize}
    \item \textbf{Pre-register hypotheses:} Reduced temptation to p-hack; forced clarity about what the study tests.
    \item \textbf{Simulator-first development:} Rapid iteration without AWS costs; enabled 350+ runs in minutes.
    \item \textbf{Manifests as source of truth:} JSON manifests enable reproducibility; analysis scripts rebuild summaries from manifests, ensuring consistency.
    \item \textbf{Incremental validation:} Unit tests → integration tests → pilot runs → full campaigns caught bugs early.
    \item \textbf{CI/CD discipline:} Automated testing on every push prevented regressions.
\end{itemize}

\section{Future Enhancements}

\begin{itemize}
    \item \textbf{Live AWS validation:} Execute subset of runs on AWS, compare metrics to simulator, quantify fidelity.
    \item \textbf{FIFO queue support:} Extend simulator to FIFO semantics (ordering, content-based deduplication), test how configuration affects FIFO reliability.
    \item \textbf{Multi-region:} Model cross-region replication, test recovery under regional failures.
    \item \textbf{Alternative brokers:} Compare SQS findings to Kafka, RabbitMQ under same fault injection.
\end{itemize}

Detailed recommendations in Chapter 6 (Conclusion).
